\setlength{\baselineskip}{20pt}
\begin{preface}
随着视觉感知系统在智能交通、公共安全与人机协同等场景中的深入应用，模型在复杂环境下的稳定性与可部署性成为研究重点。作者在前期研究中围绕行人重识别任务开展了持续探索，并进一步将研究边界扩展到分类、检测和分割等典型判别任务。研究过程中逐步形成如下认识：其一，复杂噪声与分布偏移是制约判别模型泛化能力的重要因素；其二，扩散模型具备优良的分布建模与噪声抑制能力，但直接移植到判别任务常伴随推理时延增加；其三，若能在建模阶段引入去噪能力、在部署阶段消除冗余计算，可在性能与效率之间取得更优平衡。

基于上述认识，本文尝试构建“统一建模、等价融合、能力增强”的技术路线，并形成了由基础方法到增强方法的递进式研究框架。论文写作遵循“问题导向、理论支撑、实验验证、工程可用”的组织原则，力求在方法论与应用落地之间建立清晰联系。希望本研究能够为扩散思想在判别学习中的进一步发展提供可复用的分析视角与实践经验。
\end{preface}
